% ==============================================================
%  Préambule commun aux fiches -- Advanced Probabilities (M1)
%  Usage, dans chaque fiche :
%     \documentclass[11pt,a4paper]{article}
%     \input{../preambule.tex}
%     \fiche{C5}{Vecteurs gaussiens}{Gaussian vectors}
% ==============================================================
\usepackage[T1]{fontenc}
\usepackage[utf8]{inputenc}
\usepackage[french]{babel}
\usepackage{lmodern}
\usepackage[a4paper,margin=1.9cm,top=1.7cm,bottom=1.9cm]{geometry}
\usepackage{amsmath,amssymb,amsthm,mathtools}
\usepackage[inline]{enumitem}
\usepackage{booktabs,array,tabularx}
\usepackage{xcolor}
\usepackage[most]{tcolorbox}
\usepackage{fancyhdr}
\usepackage{hyperref}
\hypersetup{colorlinks=true,linkcolor=blue!50!black,urlcolor=blue!50!black}

\setlist{itemsep=1pt,topsep=3pt,parsep=0pt}
\setlength{\parindent}{0pt}
\setlength{\parskip}{3pt}
\emergencystretch=3em

% ---------- Identification de la fiche -------------------------
\newcommand{\ficheid}{}
\newcommand{\fichetitre}{}
\newcommand{\fiche}[3]{%
  \renewcommand{\ficheid}{#1}\renewcommand{\fichetitre}{#2}%
  \pagestyle{fancy}\fancyhf{}%
  \fancyhead[L]{\small\textsf{Advanced Probabilities -- Fiche #1}}%
  \fancyhead[R]{\small\textsf{#2}}%
  \fancyfoot[C]{\small\thepage}%
  \begin{center}
    {\LARGE\bfseries Fiche #1 -- #2}\par\smallskip
    {\large\itshape (#3)}\par\smallskip
    {\small Cours de V.\,M.~Lambert, M1 -- semestre 1}
  \end{center}\vspace{-2pt}\hrule\vspace{6pt}}

% ---------- Environnements numérotés (un seul compteur) --------
\theoremstyle{definition}
\newtheorem{definition}{Définition}
\renewcommand{\thedefinition}{\ficheid.\arabic{definition}}
\newtheorem{theoreme}[definition]{Théorème}
\newtheorem{proposition}[definition]{Proposition}
\newtheorem{lemme}[definition]{Lemme}
\newtheorem{corollaire}[definition]{Corollaire}
\newtheorem{exemple}[definition]{Exemple}
\theoremstyle{remark}
\newtheorem{remarque}[definition]{Remarque}

% Résultat hors cours (standard, signalé comme tel)
\theoremstyle{definition}
\newtheorem{horscours}[definition]{Hors cours}

% Idée de preuve
\newenvironment{idee}{\par\small\textit{Idée de preuve.}~}{\par}

% ---------- Encadrés -----------------------------------------
\tcbset{fichebox/.style={enhanced,breakable,boxrule=0.6pt,arc=2pt,
  left=5pt,right=5pt,top=3pt,bottom=3pt,fonttitle=\bfseries}}
\newtcolorbox{piege}[1][]{fichebox,colback=red!4,colframe=red!60!black,
  title={Pièges classiques},#1}
\newtcolorbox{liens}[1][]{fichebox,colback=blue!3,colframe=blue!50!black,
  title={Liens entre notions},#1}
\newtcolorbox{methode}[1][Méthode]{fichebox,colback=green!4,colframe=green!40!black,title={#1}}
\newtcolorbox{alerte}[1][]{fichebox,colback=orange!6,colframe=orange!70!black,
  title={Attention à l'énoncé},#1}

% Exercice de TD (fiches de rappels)
\newcounter{exo}
\newenvironment{exo}[1][]{\refstepcounter{exo}\par\medskip
  \noindent{\large\bfseries Exercice~\theexo}\ifx\relax#1\relax\else~{\large\bfseries -- #1}\fi\par\nopagebreak\smallskip}{\par}
\newcommand{\outils}{\textbf{Outils du cours.}~}
\newcommand{\meth}{\textbf{Méthode / astuce.}~}
\newcommand{\indic}{\textbf{Indication.}~}
\newcommand{\rep}{\textbf{Réponse.}~}

% ---------- Notations du cours ---------------------------------
\newcommand{\R}{\mathbb{R}}
\newcommand{\N}{\mathbb{N}}
\newcommand{\Q}{\mathbb{Q}}
\newcommand{\C}{\mathbb{C}}
\newcommand{\PP}{\mathbb{P}}
\newcommand{\QQ}{\mathbb{Q}}
\newcommand{\E}{\mathbb{E}}
\newcommand{\F}{\mathcal{F}}
\newcommand{\G}{\mathcal{G}}
\newcommand{\Hc}{\mathcal{H}}
\newcommand{\Tc}{\mathcal{T}}
\newcommand{\Bc}{\mathcal{B}}
\newcommand{\Cc}{\mathcal{C}}
\newcommand{\Ec}{\mathcal{E}}
\newcommand{\Pc}{\mathcal{P}}
\newcommand{\Nc}{\mathcal{N}}
\newcommand{\LN}{\mathcal{LN}}
\newcommand{\Expo}{\mathcal{E}}
\newcommand{\Poi}{\mathcal{P}}
\newcommand{\Ff}{\mathbb{F}}
\newcommand{\1}{\mathbf{1}}
\newcommand{\dd}{\,\mathrm{d}}
\DeclareMathOperator{\Var}{\mathbb{V}}
\DeclareMathOperator{\Cov}{Cov}
\DeclareMathOperator{\card}{card}
\DeclareMathOperator{\diag}{diag}
\newcommand{\indep}{\perp\!\!\!\perp}
\newcommand{\ps}{\ \PP\text{-p.s.}}
\newcommand{\tendloi}{\xrightarrow[n\to+\infty]{\text{loi}}}
\newcommand{\tendps}{\xrightarrow[n\to+\infty]{\text{p.s.}}}
\newcommand{\tendproba}{\xrightarrow[n\to+\infty]{\PP}}
\newcommand{\tendLp}{\xrightarrow[n\to+\infty]{L^p}}
\newcommand{\tend}{\xrightarrow[n\to+\infty]{}}
